\clearpage\renewcommand{\refname}{参考文献}
\phantomsection\label{app:references}\addcontentsline{toc}{section}{参考文献}
\begin{thebibliography}{338}
\small\setlength{\itemsep}{3pt plus 1pt minus 1pt}
\sourcepageanchor{540}
\bibitem[1]{source1}\label{bib:source1} Agresti, A. (1990). Categorical Data Analysis. New York: Wiley.
\bibitem[2]{source2}\label{bib:source2} Agresti, A. and Caffo, B. (2000). Simple and Effective Confidence Intervals for Proportions and Differences of Proportions Result from Adding Two Successes and Two Failures. Amer. Statist. 54 280–288.
\bibitem[3]{source3}\label{bib:source3} Agresti, A. and Coull, B. A. (1998). Approximate Is Better Than “Exact” for Interval Estimation of Binomial Proportions. Amer. Statist. 52 119–126.
\bibitem[4]{source4}\label{bib:source4} Anderson, T. W. (1976). Estimation of Linear Functional Relationships: Approximate Distributions and Connection with Simultaneous Equations in Econometrics (with discussion). J. Roy. Statist. Soc. Ser. B 38 1–36.
\bibitem[5]{source5}\label{bib:source5} Anderson, T. W. (1984). An Introduction to Multivariate Statistical Analysis, 2nd edition. New York: Wiley.
\bibitem[6]{source6}\label{bib:source6} Balanda, K. P and MacGillivray, H. L. (1988). Kurtosis: A Critical Review. Amer. Statist. 42 111–119.
\bibitem[7]{source7}\label{bib:source7} Bar-Lev, S. K. and Enis, P. (1988). On the Classical Choice of Variance Stabilizing Transformations and an Application for a Poisson Variate. Biometrika 75 803–804.
\bibitem[8]{source8}\label{bib:source8} Bar-Lev, S. K. and Enis, P. (1990). On the Construction of Classes of Variance Stabilizing Transformations. Statist. Prob. Let. 10 95–100.
\bibitem[9]{source9}\label{bib:source9} Barlow, R. and Proschan, F. (1975). Statistical Theory of Life Testing. New York: Holt, Rinehart and Winston.
\bibitem[10]{source10}\label{bib:source10} Barnard, G. A. (1949). Statistical Inference (with discussion). J. Roy. Statist. Soc. Ser. B 11, 115–139.
\bibitem[11]{source11}\label{bib:source11} Barnard, G. A. (1980). Pivotal Inference and the Bayesian Controversy (with discussion). Bayesian Statistics (J. M. Bernardo, M. H. DeGroot, D. V. Lindley, and A. F. M. Smith, eds.). Valencia: University Press.
\bibitem[12]{source12}\label{bib:source12} Barr, D. R. and Zehna, P. W. (1983). Probability: Modeling Uncertainty. Reading, MA: Addison-Wesley.
\bibitem[13]{source13}\label{bib:source13} Basu, D. (1959). The Family of Ancillary Statistics. Sankhya, Ser. A 21 247–256.
\bibitem[14]{source14}\label{bib:source14} Bechhofer, R. E. (1954). A Single-Sample Multiple Decision Procedure for Ranking Means of Normal Populations with Known Variances. Ann. Math. Statist. 25 16–39.
\bibitem[15]{source15}\label{bib:source15} Behboodian, J. (1990). Examples of Uncorrelated Dependent Random Variables Using a Bivariate Mixture. Amer. Statist. 44 218.
\bibitem[16]{source16}\label{bib:source16} Berg, C. (1988). The Cube of a Normal Distribution Is Indeterminate. Ann. Prob. 16 910–913.
\bibitem[17]{source17}\label{bib:source17} Berger, J. O. (1984). The Robust Bayesian Viewpoint (with discussion). Robustness of Bayesian Analysis (J. Kadane, ed.), 63–144. Amsterdam: North-Holland.
\bibitem[18]{source18}\label{bib:source18} Berger, J. O. (1985). Statistical Decision Theory and Bayesian Analysis, 2nd edition. New York: Springer-Verlag.
\bibitem[19]{source19}\label{bib:source19} Berger, J. O. (1990). Robust Bayesian Analysis: Sensitivity to the Prior. J. Statist. Plan. Inf. 25 303–328.
\sourcepageanchor{541}
\bibitem[20]{source20}\label{bib:source20} Berger, J. O. (1994). An Overview of Robust Bayesian Analysis (with discussion). Test 3 5–124.
\bibitem[21]{source21}\label{bib:source21} Berger, J. O. and Sellke, T. (1987). Testing a Point Null Hypothesis: The Irreconcilability of P Values and Evidence (with discussion). J. Amer. Statist. Assoc. 82 112–122.
\bibitem[22]{source22}\label{bib:source22} Berger, J. O. and Wolpert, R. W. (1984). The Likelihood Principle. Institute of Mathematical Statistics Lecture Notes—Monograph Series. Hayward, CA: IMS.
\bibitem[23]{source23}\label{bib:source23} Berger, R. L. (1982). Multiparameter Hypothesis Testing and Acceptance Sampling. Technometrics 24 295–300.
\bibitem[24]{source24}\label{bib:source24} Berger, R. L. (1996). More Powerful Tests from Confidence Interval P Values. Amer. Statist. 50 314–318.
\bibitem[25]{source25}\label{bib:source25} Berger, R. L. and Boos, D. D. (1994). P Values Maximized over a Confidence Set for the Nuisance Parameter. J. Amer. Statist. Assoc. 89 1012–1016.
\bibitem[26]{source26}\label{bib:source26} Berger, R. L. and Casella, G. (1992). Deriving Generalized Means as Least Squares and Maximum Likelihood Estimates. Amer. Statist. 46 279–282.
\bibitem[27]{source27}\label{bib:source27} Berger, R. L. and Hsu, J. C. (1996). Bioequivalence Trials, Intersection-Union Tests and Equivalence Confidence Sets (with discussion). Statist. Sci. 11 283–319.
\bibitem[28]{source28}\label{bib:source28} Berkson, J. (1950). Are There Two Regressions? J. Amer. Statist. Assoc. 45 164–180.
\bibitem[29]{source29}\label{bib:source29} Bernardo, J. M and Smith, A. F. M. (1994). Bayesian Theory. New York: Wiley.
\bibitem[30]{source30}\label{bib:source30} Betteley, I. G. (1977). The Addition Law for Expectations. Amer. Statist. 31 33–35.
\bibitem[31]{source31}\label{bib:source31} Bickel, P. J. and Doksum, K. A. (1981). An Analysis of Transformations Revisited. J. Amer. Statist. Assoc. 76 296–311.
\bibitem[32]{source32}\label{bib:source32} Bickel, P. J. and Mallows, C. L. (1988). A Note on Unbiased Bayes Estimates. Amer. Statist. 42 132–134.
\bibitem[33]{source33}\label{bib:source33} Billingsley P. (1995). Probability and Measure, 3rd edition. New York: Wiley.
\bibitem[34]{source34}\label{bib:source34} Binder, D. A. (1993). Letter to the Editor about the “Exchange Paradox”. Amer. Statist. 47 160.
\bibitem[35]{source35}\label{bib:source35} Birnbaum, A. (1962). On the Foundations of Statistical Inference (with discussion). J. Amer. Statist. Assoc. 57 269–306.
\bibitem[36]{source36}\label{bib:source36} Blachman, N. M. (1996) Letter to the Editor about the “Exchange Paradox”. Amer. Statist. 50 98–99.
\bibitem[37]{source37}\label{bib:source37} Blaker, H. (2000). Confidence Curves and Improved Exact Confidence Intervals for Discrete Distributions. Can. J. Statist. 28 783–798.
\bibitem[38]{source38}\label{bib:source38} Bloch, D. A. and Moses, L. E. (1988). Nonoptimally Weighted Least Squares. Amer. Statist. 42 50–53.
\bibitem[39]{source39}\label{bib:source39} Blyth, C. R. (1986). Approximate Binomial Confidence Limits. J. Amer. Statist. Assoc. 81 843–855; correction 84 636.
\bibitem[40]{source40}\label{bib:source40} Blyth, C. R. and Still, H. A. (1983). Binomial Confidence Intervals. J. Amer. Statist. Assoc. 78 108–116.
\bibitem[41]{source41}\label{bib:source41} Boos, D. D. (1992). On Generalized Score Tests (Com: 93V47 P311-312). Amer. Statist. 46 327–333.
\bibitem[42]{source42}\label{bib:source42} Boos, D. D. and Hughes-Oliver, J. M. (1998). Applications of Basu’s Theorem. Amer. Statist. 52 218–221.
\bibitem[43]{source43}\label{bib:source43} Box, G. E. P. (1954). Some Theorems on Quadratic Forms Applied in the Study of Analysis of Variances Problems I: Effect of Inequality of Variance in the One-Way Classification. Ann. Math. Statist. 25 290–302.
\sourcepageanchor{542}
\bibitem[44]{source44}\label{bib:source44} Box, G. E. P. and Cox, D. R. (1964). An Analysis of Transformations (with discussion). J. Roy. Statist. Soc. Ser. B 26 211–252.
\bibitem[45]{source45}\label{bib:source45} Box, G. E. P. and Cox, D. R. (1982). An Analysis of Transformations Revisited, Rebutted. J. Amer. Statist. Assoc. 77 209–210.
\bibitem[46]{source46}\label{bib:source46} Box, G.E.P. and Muller, M. (1958) A Note on the Generation of Random Normal Variates. Ann. Math. Statist. 29 610–611.
\bibitem[47]{source47}\label{bib:source47} Boyles, R. A. (1983). On the convergence of the EM algorithm. J. Roy. Statist. Soc. Ser. B 45 47–50.
\bibitem[48]{source48}\label{bib:source48} Brewster, J. F. and Zidek, J. V. (1974). Improving on Equivariant Estimators. Ann. Statist. 2 21–38.
\bibitem[49]{source49}\label{bib:source49} Brown, D. and Rothery, P. (1993). Models in Biology: Mathematics, Statistics and Computing. New York: Wiley
\bibitem[50]{source50}\label{bib:source50} Brown L. D. (1986). Fundamentals of Statistical Exponential Families with Applications to Statistical Decision Theory. Institute of Mathematical Statistics Lecture Notes—Monograph Series. Hayward, CA: IMS.
\bibitem[51]{source51}\label{bib:source51} Brown L. D. (1990a). Comment on the Paper by Maatta and Casella. Statist. Sci. 5 103–106.
\bibitem[52]{source52}\label{bib:source52} Brown, L. D. (1990b). An Ancillarity Paradox Which Occurs in Multiple Linear Regression (with discussion). Ann. Statist. 18 471–492.
\bibitem[53]{source53}\label{bib:source53} Brown, L. D., Casella, G. and Hwang, J. T. G. (1995). Optimal Confidence Sets, Bioequivalence, and the Limaçon of Pascal. J. Amer. Statist. Assoc. 90 880–889.
\bibitem[54]{source54}\label{bib:source54} Brown, L. D. and Purves, R. (1973). Measurable Selections of Extrema. Ann. Statist. 1 902–912.
\bibitem[55]{source55}\label{bib:source55} Brown, P. J. and Fuller, W. A. (1991). Statistical Analysis of Measurement Error Models and Applications. Providence, RI: American Mathematical Society.
\bibitem[56]{source56}\label{bib:source56} Buehler, R. J. (1982). Some Ancillary Statistics and Their Properties (with discussion). J. Amer. Statist. Assoc. 77 581–594.
\bibitem[57]{source57}\label{bib:source57} Carlin, B.P. and Louis, T.A. (1996). Bayes and Empirical Bayes Methods for Data Analysis. London: Chapman and Hall.
\bibitem[58]{source58}\label{bib:source58} Carmer, S. G. and Walker, W. M. (1982). Baby Bear’s Dilemma: A Statistical Tale. Agronomy Journal 74 122–124.
\bibitem[59]{source59}\label{bib:source59} Carroll, R. J., Gallo, P., and Gleser, L. J. (1985). Comparison of Least Squares and Errorsin-Variables Regression, with Special Reference to Randomized Analysis of Covariance. J. Amer. Statist. Assoc. 80 929–932.
\bibitem[60]{source60}\label{bib:source60} Carroll, R. J. and Ruppert, D. (1985). Transformations: A Robust Analysis. Technometrics 27 1–12.
\bibitem[61]{source61}\label{bib:source61} Carroll, R. J. and Ruppert, D. (1988). Transformation and Weighting in Regression. London: Chapman and Hall.
\bibitem[62]{source62}\label{bib:source62} Carroll, R. J., Ruppert, D., and Stefanski, L. A. (1995). Measurement Error in Nonlinear Models. London: Chapman and Hall.
\bibitem[63]{source63}\label{bib:source63} Casella, G. (1985). An Introduction to Empirical Bayes Data Analysis. Amer. Statist. 39 83–87.
\bibitem[64]{source64}\label{bib:source64} Casella, G. (1986). Refining Binomial Confidence Intervals. Can. J. Statist. 14 113–129.
\bibitem[65]{source65}\label{bib:source65} Casella, G. (1992). Illustrating Empirical Bayes Methods. Chemolab 16 107–125.
\bibitem[66]{source66}\label{bib:source66} Casella, G. and Berger, R. L. (1987). Reconciling Bayesian and Frequentist Evidence in the One-Sided Testing Problem (with discussion). J. Amer. Statist. Assoc. 82 106–111.
\sourcepageanchor{543}
\bibitem[67]{source67}\label{bib:source67} Casella, G. and George, E. I. (1992). Explaining the Gibbs Sampler. Amer. Statist. 46 167–174.
\bibitem[68]{source68}\label{bib:source68} Casella, G. and Hwang, J. T. (1987). Employing Vague Prior Information in the Construction of Confidence Sets. J. Mult. Analysis. 21 79–104.
\bibitem[69]{source69}\label{bib:source69} Casella, G. and Robert, C. (1989). Refining Poisson Confidence Intervals. Can. J. Statist. 17 45–57.
\bibitem[70]{source70}\label{bib:source70} Casella, G. and Robert, C. P. (1996). Rao-Blackwellisation of Sampling Schemes. Biometrika 83 81–94.
\bibitem[71]{source71}\label{bib:source71} Casella, G. and Strawderman, W. E. (1980). Confidence Bands for Linear Regression with Restricted Predictor Variables. J. Amer. Statist. Assoc. 75 862–868.
\bibitem[72]{source72}\label{bib:source72} Chapman, D. G. and Robbins, H. (1951). Minimum Variance Estimation Without Regularity Assumptions. Ann. Math. Statist. 22 581–586.
\bibitem[73]{source73}\label{bib:source73} Chib, S. and Greenberg, E. (1995). Understanding the Metropolis–Hastings Algorithm. Amer. Statist. 49 327–335.
\bibitem[74]{source74}\label{bib:source74} Chikkara, R. S. and Folks, J. L. (1989). The Inverse Gaussian Distribution: Theory, Methodology, and Applications. New York: Marcel Dekker.
\bibitem[75]{source75}\label{bib:source75} Christensen, R. (1996). Plane Answers to Complex Questions. The Theory of Linear Models, 2nd edition. New York: Springer-Verlag.
\bibitem[76]{source76}\label{bib:source76} Christensen, R. and Utts, J. (1992). Bayesian Resolution of the “Exchange Paradox.” Amer. Statist. 46 274–278.
\bibitem[77]{source77}\label{bib:source77} Chun, Y. H. (1999). On the Information Economics Approach to the Generalized Game Show Problem. Amer. Statist. 53 43–51.
\bibitem[78]{source78}\label{bib:source78} Chung, K. L. (1974). A Course in Probability Theory. New York: Academic Press.
\bibitem[79]{source79}\label{bib:source79} Clopper, C. J. and Pearson, E. S. (1934). The Use of Confidence or Fiducial Limits Illustrated in the Case of the Binomial. Biometrika 26 404–413. (Also in The Selected Papers of E. S. Pearson, New York: Cambridge University Press, 1966.)
\bibitem[80]{source80}\label{bib:source80} Cochran, W. G. (1934). The Distribution of Quadratic Forms in a Normal System with Applications to the Analysis of Variance. Proceedings of the Cambridge Philosophical Society 30 178–191.
\bibitem[81]{source81}\label{bib:source81} Cochran, W. G. (1968). Errors of Measurement in Statistics. Technometrics 10 637–666.
\bibitem[82]{source82}\label{bib:source82} Cochran, W. G. and Cox, G. M. (1957). Experimental Designs, 2nd edition. New York: Wiley.
\bibitem[83]{source83}\label{bib:source83} Cornfield, J. (1969). The Bayesian Outlook and Its Application (with discussion). Biometrika 25 617–657.
\bibitem[84]{source84}\label{bib:source84} Cox, D. R. (1958). Some Problems Connected with Statistical Inference. Ann. Math. Statist. 29 357–372.
\bibitem[85]{source85}\label{bib:source85} Cox, D. R. (1971). The Choice Between Ancillary Statistics. J. Roy. Statist. Soc. Ser. B 3, 251–255.
\bibitem[86]{source86}\label{bib:source86} Cox, D. R. and Hinkley, D. V. (1974). Theoretical Statistics. London: Chapman and Hall.
\bibitem[87]{source87}\label{bib:source87} Creasy, M. A. (1956). Confidence Limits for the Gradient in the Linear Functional Relationship. J. Roy. Statist. Soc. Ser. B 18 65–69.
\bibitem[88]{source88}\label{bib:source88} Crow, E. L. (1956). Confidence Intervals for a Proportion. Biometrika 43 423–425.
\bibitem[89]{source89}\label{bib:source89} Crow, E. L. and Gardner, R. S. (1959). Confidence Intervals for the Expectation of a Poisson Variable. Biometrika 46 441–453.
\bibitem[90]{source90}\label{bib:source90} Curtiss, J. H. (1943). On Transformations Used in the Analysis of Variance. Ann. Math. Statist. 14 107–122.
\sourcepageanchor{544}
\bibitem[91]{source91}\label{bib:source91} Dalal, S. R., Fowlkes, E. B., and Hoadley, B. (1989). Risk Analysis of the Space Shuttle: Pre-Challenger Prediction of Failure. J. Amer. Statist. Assoc. 84 945–957.
\bibitem[92]{source92}\label{bib:source92} David, H. A. (1985). Bias of $S^2$ Under Dependence. Amer. Statist. 39 201.
\bibitem[93]{source93}\label{bib:source93} Davidson, R. R. and Solomon, D. L. (1974). Moment-Type Estimation in the Exponential Family. Communications in Statistics 3 1101–1108.
\bibitem[94]{source94}\label{bib:source94} Dean, A. and Voss, D. (1999) Design and Analysis of Experiments New York: Springer- Verlag.
\bibitem[95]{source95}\label{bib:source95} deFinetti, B. (1972). Probability, Induction, and Statistics. New York: Wiley.
\bibitem[96]{source96}\label{bib:source96} DeGroot, M. H. (1986). Probability and Statistics, 2nd edition. New York: Addison-Wesley.
\bibitem[97]{source97}\label{bib:source97} Dempster, A. P., Laird, N. M., and Rubin, D. B. (1977). Maximum Likelihood from Incomplete Data via the EM Algorithm. J. Roy. Statist. Soc. Ser. B 39 1–22.
\bibitem[98]{source98}\label{bib:source98} Devroye, L. (1985). Non-Uniform Random Variate Generation. New York: Springer-Verlag.
\bibitem[99]{source99}\label{bib:source99} Diaconis, P. and Holmes, S. (1994). Gray Codes for Randomization Procedures. Statistics and Computing 4 287–302.
\bibitem[100]{source100}\label{bib:source100} Diaconis, P. and Mosteller, F. (1989). Methods for Studying Coincidences. J. Amer. Statist. Assoc. 84 853–861.
\bibitem[101]{source101}\label{bib:source101} Draper, D. (1988). Rank-Based Robust Analysis of Linear Models. I. Exposition and Review (with discussion). Statist. Sci. 3 239–271.
\bibitem[102]{source102}\label{bib:source102} Draper, N. R. and Smith, H. (1998). Applied Regression Analysis, 3rd edition. New York: Wiley.
\bibitem[103]{source103}\label{bib:source103} Durbin, J. (1970). On Birnbaum’s Theorem and the Relation Between Sufficiency, Conditionality, and Likelihood. J. Amer. Statist. Assoc. 65 395–398.
\bibitem[104]{source104}\label{bib:source104} Dynkin, E. B. (1951). Necessary and Sufficient Statistics for a Family of Probability Distributions. English translation in Selected Translations in Mathematical Statistics and Probability 1 (1961), 23–41.
\bibitem[105]{source105}\label{bib:source105} Eberhardt, K. R. and Fligner, M. A. (1977). A Comparison of Two Tests for Equality of Two Proportions. Amer. Statist. 31 151–155.
\bibitem[106]{source106}\label{bib:source106} Efron, B. F. (1979a). Bootstrap Methods: Another Look at the Jackknife. Ann. Statist. 7 1–26.
\bibitem[107]{source107}\label{bib:source107} Efron, B. F. (1979b). Computers and the Theory of Statistics: Thinking the Unthinkable. SIAM Review 21 460–480.
\bibitem[108]{source108}\label{bib:source108} Efron, B. F. (1982). The Jackknife, the Bootstrap, and Other Resampling Plans. Philadelphia: Society for Industrial and Applied Mathematics.
\bibitem[109]{source109}\label{bib:source109} Efron, B. F. (1998). R. A. Fisher in the 21st Century (with discussion). Statist. Sci. 13 95–122.
\bibitem[110]{source110}\label{bib:source110} Efron, B. F. and Hinkley, D. V. (1978). Assessing the Accuracy of the Maximum Likelihood Estimator: Observed Versus Expected Fisher Information. Biometrika 65 457–487.
\bibitem[111]{source111}\label{bib:source111} Efron, B. F. and Morris, C. N. (1972). Limiting the Risk of Bayes and Empirical Bayes Estimators Part II: The Empirical Bayes Case. J. Amer. Statist. Assoc. 67 130–139.
\bibitem[112]{source112}\label{bib:source112} Efron, B. F. and Morris, C. N. (1973). Stein’s Estimation Rule and Its Competitors—An Empirical Bayes Approach. J. Amer. Statist. Assoc. 68 117–130.
\bibitem[113]{source113}\label{bib:source113} Efron, B. F. and Morris, C. N. (1975). Data Analysis Using Stein’s Estimator and Its Generalizations. J. Amer. Statist. Assoc. 70 311–319.
\bibitem[114]{source114}\label{bib:source114} Efron, B. and Tibshirani, R. J. (1993). An Introduction to the Bootstrap. London: Chapman and Hall..
\bibitem[115]{source115}\label{bib:source115} Ellis, S. Instability of Least Squares, Least Absolute Deviation and Least Median of Squares Linear Regression (with discussion). Statist. Sci. 13 337–350.
\sourcepageanchor{545}
\bibitem[116]{source116}\label{bib:source116} Feldman, D. and Fox, M. (1968). Estimation of the Parameter n in the Binomial Distribution. J. Amer. Statist. Assoc. 63 150–158.
\bibitem[117]{source117}\label{bib:source117} Feldstein, M. (1974). Errors in Variables: A Consistent Estimator with Smaller MSE in Finite Samples. J. Amer. Statist. Assoc. 69 990–996.
\bibitem[118]{source118}\label{bib:source118} Feller, W. (1968). An Introduction to Probability Theory and Its Applications, Volume I. New York: Wiley.
\bibitem[119]{source119}\label{bib:source119} Feller, W. (1971). An Introduction to Probability Theory and Its Applications, Volume II. New York: Wiley.
\bibitem[120]{source120}\label{bib:source120} Ferguson, T. S. (1996). A Course in Large Sample Theory. London: Chapman and Hall.
\bibitem[121]{source121}\label{bib:source121} Fieller, E. C. (1954). Some Problems in Interval Estimation. J. Roy. Statist. Soc. Ser. B 16 175–185.
\bibitem[122]{source122}\label{bib:source122} Finch, S. J., Mendell, N. R., and Thode, H. C. (1989). Probabilistic Measures of Adequacy of a Numerical Search for a Global Maximum. J. Amer. Statist. Assoc. 84 1020–1023.
\bibitem[123]{source123}\label{bib:source123} Fisher, R. A. (1925). Theory of Statistical Estimation. Proceedings of the Cambridge Philosophical Society 22 700–725.
\bibitem[124]{source124}\label{bib:source124} Fisher, R. A. (1930). Inverse Probability. Proceedings of the Cambridge Philosophical Society 26 528–535.
\bibitem[125]{source125}\label{bib:source125} Fisher, R. A. (1935). The Fiducial Argument in Statistical Inference. Annals of Eugenics 6 391–398. (Also in R. A. Fisher, Contributions to Mathematical Statistics, New York: Wiley, 1950.)
\bibitem[126]{source126}\label{bib:source126} Fisher, R. A. (1939). The Comparison of Samples with Possibly Unequal Variances. Annals of Eugenics 9 174–180. (Also in R. A. Fisher, Contributions to Mathematical Statistics, New York: Wiley, 1950.)
\bibitem[127]{source127}\label{bib:source127} Fraser, D. A. S. (1968). The Structure of Inference. New York: Wiley.
\bibitem[128]{source128}\label{bib:source128} Fraser, D. A. S. (1979). Inference and Linear Models. New York: McGraw-Hill.
\bibitem[129]{source129}\label{bib:source129} Freedman, D., Pisani, R., Purves, R., and Adhikari, A. (1991). Statistics, 2nd edition. New York: Norton.
\bibitem[130]{source130}\label{bib:source130} Fuller, W. A. (1987). Measurement Error Models. New York: Wiley.
\bibitem[131]{source131}\label{bib:source131} Gafarian, A. V. (1964). Confidence Bands in Straight Line Regression. J. Amer. Statist. Assoc. 59 182–213.
\bibitem[132]{source132}\label{bib:source132} Gardner, M. (1961). The Second Scientific American Book of Mathematical Puzzles and Diversions. New York: Simon \& Schuster.
\bibitem[133]{source133}\label{bib:source133} Garwood, F. (1936). Fiducial Limits for the Poisson Distribution. Biometrika 28 437–442.
\bibitem[134]{source134}\label{bib:source134} Gelfand, A. E. and Smith, A. F. M. (1990). Sampling-Based Approaches to Calculating Marginal Densities. J. Amer. Statist. Assoc. 85 398–409.
\bibitem[135]{source135}\label{bib:source135} Gelman, A. and Meng, X-L. (1991). A Note on Bivariate Distributions That Are Conditionally Normal. Amer. Statist. 45 125–126.
\bibitem[136]{source136}\label{bib:source136} Gelman, A. and Rubin, D.B. (1992). Inference from Iterative Simulation Using Multiple Sequences (with discussion). Statist. Sci. 7 457–511.
\bibitem[137]{source137}\label{bib:source137} Geman, S. and Geman, D. (1984) Stochastic Relaxation, Gibbs Distributions and the Bayesian Restoration of Images. IEEE Trans. Pattern Anal. Mach. Intell. 6 721–741.
\bibitem[138]{source138}\label{bib:source138} Geyer, C.J. and Thompson, E.A. (1992) Constrained Monte Carlo Maximum Likelihood for Dependent Data (with discussion). J. Roy. Statist. Soc. Ser. B 54 657–699.
\bibitem[139]{source139}\label{bib:source139} Ghosh, B. K. (1979). A Comparison of Some Approximate Confidence Intervals for the Binomial Parameter. J. Amer. Statist. Assoc. 74 894–900.
\sourcepageanchor{546}
\bibitem[140]{source140}\label{bib:source140} Ghosh, M. and Meeden, G. (1977). On the Non-Attainability of Chebychev Bounds. Amer. Statist. 31 35–36.
\bibitem[141]{source141}\label{bib:source141} Gianola, D. and Fernando, R. L. (1986). Bayesian Methods in Animal Breeding Theory. Journal of Animal Science 63 217–244.
\bibitem[142]{source142}\label{bib:source142} Gilat, D. (1977). Monotonicity of a Power Function: An Elementary Probabilistic Proof. Amer. Statist. 31 91–93.
\bibitem[143]{source143}\label{bib:source143} Gleser, L. J. (1981). Estimation in a Multivariate Errors-in-Variables Regression Model: Large-Sample Results. Ann. Statist. 9 24–44.
\bibitem[144]{source144}\label{bib:source144} Gleser, L. J. (1983). Functional, Structural, and Ultrastructural Errors-in-Variables Models. Proceedings of the Business and Economic Statistics Section 57–66. Alexandria, VA: American Statistical Association.
\bibitem[145]{source145}\label{bib:source145} Gleser, L. J. (1987). Confidence Intervals for the Slope in a Linear Errors-in-Variables Regression Model. Advances in Multivariate Statistical Analysis (K. Gupta, ed.) 85–109. Dordrecht: D. Reidel.
\bibitem[146]{source146}\label{bib:source146} Gleser, L. J. (1989). The Gamma Distribution as a Mixture of Exponential Distributions. Amer. Statist. 43 115–117.
\bibitem[147]{source147}\label{bib:source147} Gleser, L. J. (1991). Measurement error models (with discussion). Chem. Int. Lab. Sys. 10 45–67.
\bibitem[148]{source148}\label{bib:source148} Gleser, L. J., Carroll, R. J., and Gallo, P. (1987). The Limiting Distribution of Least Squares in an Errors-in-Variables Regression Model. Ann. Statist. 15 220–233.
\bibitem[149]{source149}\label{bib:source149} Gleser, L. J. and Healy, J. D. (1976). Estimating the Mean of a Normal Distribution with Known Coefficient of Variation. J. Amer. Statist. Assoc. 71 977–981.
\bibitem[150]{source150}\label{bib:source150} Gleser, L. J. and Hwang, J. T. (1987). The Nonexistence of 100(1 − \ensuremath{\alpha})\% Confidence Sets of Finite Expected Diameter in Errors-in-Variables and Related Models. Ann. Statist. 15 1351–1362.
\bibitem[151]{source151}\label{bib:source151} Gnedenko, B. V. (1978). The Theory of Probability. Moscow: MIR Publishers.
\bibitem[152]{source152}\label{bib:source152} Groeneveld, R. A. (1991). An Influence Function Approach to Describing the Skewness of a Distribution. Amer. Statist. 45 97–102.
\bibitem[153]{source153}\label{bib:source153} Guenther, W. C. (1978). Some Easily Found Minimum Variance Unbiased Estimators. Amer. Statist. 32 29–33.
\bibitem[154]{source154}\label{bib:source154} Hall, P. (1992). The Bootstrap and Edgeworth Expansion. New York: Springer-Verlag.
\bibitem[155]{source155}\label{bib:source155} Halmos, P. R. and Savage, L. J. (1949). Applications of the Radon–Nikodym Theorem to the Theory of Sufficient Statistics. Ann. Math. Statist. 20 225–241.
\bibitem[156]{source156}\label{bib:source156} Hampel, F. R. (1974). The Influence Curve and Its Role in Robust Estimation. J. Amer. Statist. Assoc. 69 383–393.
\bibitem[157]{source157}\label{bib:source157} Hanley, J. A. (1992). Jumping to Coincidences: Defying Odds in the Realm of the Preposterous. Amer. Statist. 46 197–202.
\bibitem[158]{source158}\label{bib:source158} Hardy, G. H., Littlewood, J. E., and Polya, G. (1952). Inequalities, 2nd edition. London: Cambridge University Press.
\bibitem[159]{source159}\label{bib:source159} Hartley, H. O. (1958). Maximum Likelihood Estimation from Incomplete Data. Biometrics 14 174–194.
\bibitem[160]{source160}\label{bib:source160} Harville, D. A. (1981). Unbiased and Minimum-Variance Unbiased Estimation of Estimable Functions for Fixed Linear Models with Arbitrary Covariance Structure. Ann. Statist. 9 633–637.
\bibitem[161]{source161}\label{bib:source161} Hawkins, D. M. (1993). The Feasible Set Algorithm for Least Median of Squares Regression. Computational Statistics and Data Analysis 16 81–101.
\sourcepageanchor{547}
\bibitem[162]{source162}\label{bib:source162} Hawkins, D. M. (1994). The Feasible Solution Algorithm for Least Trimmed Squares Regression. Computational Statistics and Data Analysis 17 186–196.
\bibitem[163]{source163}\label{bib:source163} Hawkins, D. M. (1995). Convergence of the Feasible Solution Algorithm for Least Median of Squares Regressionn. Computational Statistics and Data Analysis 19 519–538.
\bibitem[164]{source164}\label{bib:source164} Hayter, A. J. (1984). A Proof of the Conjecture that the Tukey-Kramer Multiple Comparison Procedure Is Conservative. Ann. Statist. 12 61–75.
\bibitem[165]{source165}\label{bib:source165} Hettmansperger, T. P. and McKean, J. W. (1998). Robust Nonparametric Statistical Methods. London: Kendall’s Library of Statistics, 5, Edward Arnold; New York: Wiley.
\bibitem[166]{source166}\label{bib:source166} Hinkley, D. V. (1980). Likelihood. Can. J. Statist. 8 151–163.
\bibitem[167]{source167}\label{bib:source167} Hinkley, D. V., Reid, N, and Snell, L. (1991). Statistical Theory and Modelling. In honor of Sir David Cox. London: Chapman and Hall.
\bibitem[168]{source168}\label{bib:source168} Hinkley, D. V. and Runger, G. (1984). The Analysis of Transformed Data (with discussion). J. Amer. Statist. Assoc. 79 302–320.
\bibitem[169]{source169}\label{bib:source169} Hsu, J. C. (1996). Multiple Comparisons: Theory and Methods. London: Chapman and Hall..
\bibitem[170]{source170}\label{bib:source170} Huber, P. J. (1964). Robust Estimation of a Location Parameter. Ann. Math. Statist. 35 73–101.
\bibitem[171]{source171}\label{bib:source171} Huber, P. J. (1981). Robust Statistics. New York: Wiley.
\bibitem[172]{source172}\label{bib:source172} Hudson, H. M. (1978). A Natural Identity for Exponential Families with Applications in Multiparameter Estimation. Ann. Statist. 6 473–484.
\bibitem[173]{source173}\label{bib:source173} Huzurbazar, V. S. (1949). On a Property of Distributions Admitting Sufficient Statistics. Biometrika 36 71–74.
\bibitem[174]{source174}\label{bib:source174} Hwang, J. T. (1982). Improving on Standard Estimators in Discrete Exponential Families with Applications to Poisson and Negative Binomial Cases. Ann. Statist. 10 857–867.
\bibitem[175]{source175}\label{bib:source175} Hwang, J. T. (1995). Fieller’s Problems and Resampling Techniques. Statistica Sinica 5 161–171.
\bibitem[176]{source176}\label{bib:source176} James, W. and Stein, C. (1961). Estimation with Quadratic Loss. Proc. Fourth Berkeley Symp. Math. Statist. Prob. 1. University of California Press.
\bibitem[177]{source177}\label{bib:source177} Johnson, N. L. and Kotz, S. (1969-1972). Distributions in Statistics (4 vols.). New York: Wiley.
\bibitem[178]{source178}\label{bib:source178} Johnson, N. L., Kotz. S., and Balakrishnan, N. (1994). Continuous Univariate Distributions, Volume 1, 2nd edition. New York: Wiley.
\bibitem[179]{source179}\label{bib:source179} Johnson, N. L., Kotz. S., and Balakrishnan, N. (1995). Continuous Univariate Distributions, Volume 2, 2nd edition. New York: Wiley.
\bibitem[180]{source180}\label{bib:source180} Johnson, N. L., Kotz. S., and Kemp, A. W. (1992). Univariate Discrete Distributions, 2nd edition. New York: Wiley.
\bibitem[181]{source181}\label{bib:source181} Jones, M. C. (1999). Distributional Relationships Arising from Simple Trigonometric Formulas. Amer. Statist. 53 99–102
\bibitem[182]{source182}\label{bib:source182} Joshi, S. M. and Nabar, S. P. (1989). Linear Estimators for the Parameter in the Problem of the Nile. Amer. Statist. 43 40–41.
\bibitem[183]{source183}\label{bib:source183} Joshi, V. M. (1969). Admissibility of the Usual Confidence Sets for the Mean of a Univariate or Bivariate Normal Population. Ann. Math. Statist. 40 1042–1067.
\bibitem[184]{source184}\label{bib:source184} Juola, R. C. (1993). More on Shortest Confidence Intervals. Amer. Statist. 47 117–119.
\bibitem[185]{source185}\label{bib:source185} Kalbfleisch, J. D. (1975). Sufficiency and Conditionality. Biometrika 62 251–268.
\bibitem[186]{source186}\label{bib:source186} Kalbfleisch, J. D. and Prentice, R. L. (1980). The Statistical Analysis of Failure Time Data. New York: Wiley.
\sourcepageanchor{548}
\bibitem[187]{source187}\label{bib:source187} Karlin, S. and Ost, F. (1988). Maximal Length of Common Words among Random Letter Sequences. Ann. Prob. 16 535–563.
\bibitem[188]{source188}\label{bib:source188} Kelker, D. (1970). Distribution Theory of Spherical Distributions and a Location–Scale Parameter Generalization. Sankhyā, Ser. A 32 419–430.
\bibitem[189]{source189}\label{bib:source189} Kendall, M. and Stuart, A. (1979). The Advanced Theory of Statistics, Volume II: Inference and Relationship, 4th edition. New York: Macmillan.
\bibitem[190]{source190}\label{bib:source190} Kirk, R. E. (1982). Experimental Design: Procedures for the Behavorial Sciences, 2nd edition. Pacific Grove, CA: Brooks/Cole.
\bibitem[191]{source191}\label{bib:source191} Koopmans, L. H. (1993). A Note on Using the Moment Generating Function to Teach the Laws of Large Numbers. Amer. Statist. 47 199–202.
\bibitem[192]{source192}\label{bib:source192} Kuehl, R. O. (2000). Design of Experiments : Statistical Principles of Research Design and Analysis, 2nd edition. Pacific Grove, CA : Duxbury.
\bibitem[193]{source193}\label{bib:source193} Lange, N., Billard, L., Conquest, L., Ryan, L., Brillinger, D., Greenhouse, J. (eds). (1994), Case Studies in Biometry. New York: Wiley-Interscience.
\bibitem[194]{source194}\label{bib:source194} Le Cam, L. (1953). On Some Asymptotic Properties of Maximum Likelihood Estimates and Related Bayes’ Estimates. Univ. of Calif. Publ. in Statist. 1 277–330.
\bibitem[195]{source195}\label{bib:source195} Leemis, L. M. (1986). Relationships Among Common Univariate Distributions. Amer. Statist. 40 143–146.
\bibitem[196]{source196}\label{bib:source196} Leemis, L. M. and Trivedi, K. S. (1996). A Comparison of Approximate Interval Estimators for the Bernoulli Parameter. Amer. Statist. 50 63–68.
\bibitem[197]{source197}\label{bib:source197} Lehmann, E. L. (1981). An Interpretation of Completeness and Basu’s Theorem. J. Amer. Statist. Assoc. 76 335–340.
\bibitem[198]{source198}\label{bib:source198} Lehmann, E. L. (1986). Testing Statistical Hypotheses, 2nd edition. New York: Wiley.
\bibitem[199]{source199}\label{bib:source199} Lehmann, E. L. (1999). Introduction to Large-Sample Theory. New York: Springer-Verlag.
\bibitem[200]{source200}\label{bib:source200} Lehmann, E. L. and Casella, G. (1998). Theory of Point Estimation, 2nd edition. New York: Springer-Verlag.
\bibitem[201]{source201}\label{bib:source201} Lehmann, E. L. and Scheffé, H. (1950, 1955, 1956). Completeness, Similar Regions, and Unbiased Estimation. Sankhyā, Ser. A 10 305–340; 15 219–236; correction 17 250.
\bibitem[202]{source202}\label{bib:source202} Lehmann, E. L. and Scholz, F. W. (1992). Ancillarity. Current Issues in Statistical Inference: Essays in Honor of D. Basu (M. Ghosh and P. K. Pathak, eds.). Hayward, CA: IMS Monograph Series, 32–51.
\bibitem[203]{source203}\label{bib:source203} LePage, R. and Billard, L. (eds.) (1992). Exploring the Limits of Bootstrap. New York: Wiley.
\bibitem[204]{source204}\label{bib:source204} Lindley, D. V. (1957). A Statistical Paradox. Biometrika 44 187–192.
\bibitem[205]{source205}\label{bib:source205} Lindley, D. V. (1962). Discussion of the Article by Stein. J. Roy. Statist. Soc. Ser. B 24 265–296.
\bibitem[206]{source206}\label{bib:source206} Lindley, D. V. and Phillips, L. D. (1976). Inference for a Bernoulli Process (a Bayesian View). Amer. Statist. 30 112–119.
\bibitem[207]{source207}\label{bib:source207} Lindley, D. V. and Smith, A. F. M. (1972). Bayes Estimates for the Linear Model. J. Roy. Statist. Soc. Ser. B 34 1–41.
\bibitem[208]{source208}\label{bib:source208} Little, R. J. A. and Rubin, D. B. (1987). Statistical Analysis with Missing Data. New York: Wiley.
\bibitem[209]{source209}\label{bib:source209} Liu, R. Y. (1990). On a Notion of Data Depth Based on Random Simplices. Ann. Statist. 18 405–414.
\bibitem[210]{source210}\label{bib:source210} Liu, R. Y. and Singh, K. (1992). Ordering Directional Data: Concepts of Data Depth on Circles and Spheres. Ann. Statist. 20 1468–1484.
\sourcepageanchor{549}
\bibitem[211]{source211}\label{bib:source211} Luceño, A. (1997). Further Evidence Supporting the Numerical Usefulness of Characteristic Functions. Amer. Statist. 51 233–234.
\bibitem[212]{source212}\label{bib:source212} Maatta, J. M. and Casella, G. (1987). Conditional Properties of Interval Estimators of the Normal Variance. Ann. Statist. 15 1372–1388.
\bibitem[213]{source213}\label{bib:source213} Madansky, A. (1962). More on Length of Confidence Intervals. J. Amer. Statist. Assoc. 57 586–589.
\bibitem[214]{source214}\label{bib:source214} Marshall, A. W. and Olkin, I. (1979). Inequalities: Theory of Majorization and Its Applications. New York: Academic Press.
\bibitem[215]{source215}\label{bib:source215} McCullagh, P. (1994). Does the Moment Generating Function Characterize a Distribution? Amer. Statist. 48 208.
\bibitem[216]{source216}\label{bib:source216} McLachlan, G., and Krishnan, T. (1997). The EM Algotithm and Extensions. New York: Wiley.
\bibitem[217]{source217}\label{bib:source217} McPherson, G. (1990). Statistics in Scientific Inverstigation. New York: Springer-Verlag.
\bibitem[218]{source218}\label{bib:source218} Mengersen, K. L. and Tweedie, R. L. (1996). Rates of Convergence of the Hastings and Metropolis Algorithms. Ann. Statist. 24 101–121.
\bibitem[219]{source219}\label{bib:source219} Metropolis, N., Rosenbluth, A. W., Rosenbluth, M. N., Teller, A. H., Teller, E. (1953). Equations of State Calculations by Fast Computing Machines. J. Chem. Phys. 21 1087–1092.
\bibitem[220]{source220}\label{bib:source220} Miller, R. G. (1974). The Jackknife—A Review. Biometrika 61 1–15.
\bibitem[221]{source221}\label{bib:source221} Miller, R. G. (1981). Simultaneous Statistical Inference, 2nd edition. New York: Springer- Verlag.
\bibitem[222]{source222}\label{bib:source222} Moran, P. A. P. (1971). Estimating Structural and Functional Relationships. J. Mult. Analysis. 1 232–255.
\bibitem[223]{source223}\label{bib:source223} Morgan, J. P., Chaganty, N. R., Dahiya, R. C. and Doviak, M. J. (1991). Let’s Make a Deal: The Player’s Dilemma (with discussion). Amer. Statist. 45 284–289.
\bibitem[224]{source224}\label{bib:source224} Morris, C. N. (1982). Natural Exponential Families with Quadratic Variance Functions. Ann. Statist. 10 65–80.
\bibitem[225]{source225}\label{bib:source225} Morris, C. N. (1983). Parametric Empirical Bayes Inference: Theory and Applications (with discussion). J. Amer. Statist. Assoc. 78 47–65.
\bibitem[226]{source226}\label{bib:source226} Morrison, D. G. (1978). A Probability Model for Forced Binary Choices. Amer. Statist. 32 23–25.
\bibitem[227]{source227}\label{bib:source227} Naiman, D. Q. (1983). Comparing Scheffé-Type to Constant Width Confidence Bounds in Regression. J. Amer. Statist. Assoc. 78 906–912.
\bibitem[228]{source228}\label{bib:source228} Naiman, D. Q. (1984). Optimal Simultaneous Confidence Bounds. Ann. Statist. 12 702–715.
\bibitem[229]{source229}\label{bib:source229} Naiman, D. Q. (1987). Simultaneous Confidence Bounds in Multiple Regression Using Predictor Variable Constraints. J. Amer. Statist. Assoc. 82 214–219.
\bibitem[230]{source230}\label{bib:source230} Naiman, D. Q. and Wynn, H. P. (1992). Inclusion-Exclusion-Bonferroni Identities and Inequalities for Discrete Tube-like Problems via Euler Characteristics. Ann. Statist. 20 43–76.
\bibitem[231]{source231}\label{bib:source231} Naiman, D. Q. and Wynn, H. P. (1997). Abstract Tubes, Improved Inclusion-Exclusion Identities and Inequalities and Importance Sampling. Ann. Statist. 25 1954–1983.
\bibitem[232]{source232}\label{bib:source232} Neter, J., Wasserman, W., and Whitmore, G. A. (1993). Applied Statistics. Boston : Allyn and Bacon.
\bibitem[233]{source233}\label{bib:source233} Neyman, J. (1935). Su un teorema concernente le cosiddette statistiche sufficienti. Inst. Ital. Atti. Giorn. 6 320–334.
\bibitem[234]{source234}\label{bib:source234} Noorbaloochi, S. and Meeden, G. (1983). Unbiasedness as the Dual of Being Bayes. J. Amer. Statist. Assoc. 78 619–623.
\sourcepageanchor{550}
\bibitem[235]{source235}\label{bib:source235} Norton, R. M. (1984). The Double Exponential Distribution: Using Calculus to Find an MLE. Amer. Statist. 38 135–136.
\bibitem[236]{source236}\label{bib:source236} Nussbaum, M. (1976). Maximum Likelihood and Least Squares Estimation of Linear Functional Relationships. Mathematische Operationsforschung und Statistik, Ser. Statistik 7 23–49.
\bibitem[237]{source237}\label{bib:source237} Olkin, I., Petkau, A. J., and Zidek, J. V. (1981). A Comparison of n Estimators for the Binomial Distribution. J. Amer. Statist. Assoc. 76 637–642.
\bibitem[238]{source238}\label{bib:source238} Pal, N. and Berry, J. (1992). On Invariance and Maximum Likelihood Estimation. Amer. Statist. 46 209–212.
\bibitem[239]{source239}\label{bib:source239} Park, C. G., Park, T., and Shin, D. W. (1996). A Simple Method for Generating Correlated Binary Variates. Amer. Statist. 50 306–310.
\bibitem[240]{source240}\label{bib:source240} Pena, E. A. and Rohatgi, V. (1994). Some Comments about Sufficiency and Unbiased Estimation. Amer. Statist. 48 242–243.
\bibitem[241]{source241}\label{bib:source241} Piegorsch, W. W. (1985). Admissible and Optimal Confidence Bands in Simple Linear Regression. Ann. Statist. 13 801–810.
\bibitem[242]{source242}\label{bib:source242} Pitman, E. J. G. (1939). The Estimation of the Location and Scale Parameters of a Continuous Population of Any Given Form. Biometrika 30 200–215.
\bibitem[243]{source243}\label{bib:source243} Portnoy, S. (1987). A Central Limit Theorem Applicable to Robust Regression Estimators. J. Mult. Analysis. 22 24–50.
\bibitem[244]{source244}\label{bib:source244} Portnoy, S. and Koenker, R. (1997). The Gaussian Hare and the Laplacian Tortoise: Computability of Squared-Error versus Absolute-Error Estimators. Statist. Sci. 12 279–300.
\bibitem[245]{source245}\label{bib:source245} Portnoy, S. and Mizera, I. (1998). Discussion of the Paper by Ellis. Statist. Sci. 13 344–347.
\bibitem[246]{source246}\label{bib:source246} Pratt, J. W. (1961). Length of Confidence Intervals. J. Amer. Statist. Assoc. 56 549–567.
\bibitem[247]{source247}\label{bib:source247} Proschan, M. A. and Presnell, B. (1998). Expect the Unexpected from Conditional Expectation. Amer. Statist. 52 248–252.
\bibitem[248]{source248}\label{bib:source248} Pukelsheim, F. (1994). The Three Sigma Rule. Amer. Statist. 48 88–91.
\bibitem[249]{source249}\label{bib:source249} Quenouille, M. H. (1956). Notes on Bias in Estimation. Biometrika 43 353–360.
\bibitem[250]{source250}\label{bib:source250} Reid, N. (1995). The Role of Conditioning in Inference (with discussion). Statist. Sci. 10 138–166.
\bibitem[251]{source251}\label{bib:source251} Resnick, S. I. (1999). A Probability Path. Basel: Birkhauser.
\bibitem[252]{source252}\label{bib:source252} Ridgeway, T. (1993). Letter to the Editor about the “Exchange Paradox.” Amer. Statist. 47 311.
\bibitem[253]{source253}\label{bib:source253} Ripley, B. D. (1987) Stochastic Simulation. New York: Wiley.
\bibitem[254]{source254}\label{bib:source254} Robbins, H. (1977). A Fundamental Question of Practical Statistics (letter to the editor). Amer. Statist. 31 97.
\bibitem[255]{source255}\label{bib:source255} Robert, C. P. (1994) The Bayesian Choice. New York: Springer-Verlag.
\bibitem[256]{source256}\label{bib:source256} Robert, C. P. and Casella, G. (1999). Monte Carlo Statistical Methods New York: Springer- Verlag.
\bibitem[257]{source257}\label{bib:source257} Robson, D. S. (1959). A Simple Method for Constructing Orthogonal Polynomials when the Independent Variable Is Unequally Spaced. Biometrics 15 187–191.
\bibitem[258]{source258}\label{bib:source258} Romano, J. P. and Siegel, A. F. (1986). Counterexamples in Probability and Statistics. Pacific Grove, CA: Wadsworth and Brooks/Cole.
\bibitem[259]{source259}\label{bib:source259} Ross, S. M. (1988) A First Course in Probability Theory, 3rd edition. New York: MacMillan.
\bibitem[260]{source260}\label{bib:source260} Ross, S. M. (1994) Letter to the Editor about the “Exchange Paradox.” Amer. Statist. 48 267.
\sourcepageanchor{551}
\bibitem[261]{source261}\label{bib:source261} Ross, S. M. (1996). Bayesians Should Not Resample a Prior Sample to Learn about the Posterior. Amer. Statist. 50 116.
\bibitem[262]{source262}\label{bib:source262} Rousseeuw, P. J. (1984). Least Median of Squares Regression. J. Amer. Statist. Assoc. 79 871–880.
\bibitem[263]{source263}\label{bib:source263} Rousseeuw, P. J. and Hubert, M. (1999). Regression Depth (with discussion). J. Amer. Statist. Assoc. 94 388–433
\bibitem[264]{source264}\label{bib:source264} Rouseeuw, P. J. and Leroy, A. M. (1987). Robust Regression and Outlier Detection. New York: Wiley.
\bibitem[265]{source265}\label{bib:source265} Royall, R. M. (1997). Statistical Evidence: A Likelihood Paradigm. London: Chapman and Hall..
\bibitem[266]{source266}\label{bib:source266} Rubin, D. B. (1988). Using the SIR Algorithm to Simulate Posterior Distributions. Bayesian Statistics 3 (J. M. Bernardo, M. H. DeGroot, D. V. Lindley, and A. F. M. Smith, eds.) Cambridge, MA: Oxford University Press 395–402.
\bibitem[267]{source267}\label{bib:source267} Rudin, W. (1976). Principles of Real Analysis. New York: McGraw-Hill.
\bibitem[268]{source268}\label{bib:source268} Ruppert, D. (1987). What Is Kurtosis? Amer. Statist. 41 1–5.
\bibitem[269]{source269}\label{bib:source269} Ruppert, D. and Carroll, R. J. (1979). Trimmed Least Squares Estimation in the Linear Model. J. Amer. Statist. Assoc. 75 828–838.
\bibitem[270]{source270}\label{bib:source270} Russell, K. G. (1991). Estimating the value of e by Simulation. Amer. Statist. 45 66–68
\bibitem[271]{source271}\label{bib:source271} Samuels, M. L. and Lu, T.-F. C. (1992). Sample Size Requirements for the Back-of-the- Envelope Binomial Confidence Interval. Amer. Statist. 46 228–231.
\bibitem[272]{source272}\label{bib:source272} Satterthwaite, F. E. (1946). An Approximate Distribution of Estimates of Variance Components. Biometrics Bulletin (now called Biometrics) 2 110–114.
\bibitem[273]{source273}\label{bib:source273} Saw, J. G., Yang, M. C. K., and Mo, T. C. (1984). Chebychev’s Inequality with Estimated Mean and Variance. Amer. Statist. 38 130–132.
\bibitem[274]{source274}\label{bib:source274} Schafer, J. L. (1997). Analysis of Incomplete Multivariate Data. London: Chapman and Hall.
\bibitem[275]{source275}\label{bib:source275} Scheffé, H. (1959). The Analysis of Variance. New York: Wiley.
\bibitem[276]{source276}\label{bib:source276} Schervish, M. J. (1995). Theory of Statistics. New York: Springer-Verlag.
\bibitem[277]{source277}\label{bib:source277} Shervisch, M. J. (1996). P Values: What They Are and What They Are Not. Amer. Statist. 50 203–206.
\bibitem[278]{source278}\label{bib:source278} Schuirmann, D. J. (1987). A Comparison of the Two One-Sided Test Procedures and the Power Approach for Assessing the Equivalence of Average Bioavailability. J. Pharmacokinetics and Biopharmaceutics 15 657–680.
\bibitem[279]{source279}\label{bib:source279} Schwager, S. J. (1984). Bonferroni Sometimes Loses. Amer. Statist. 38 192–197.
\bibitem[280]{source280}\label{bib:source280} Schwager, S. J. (1985). Reply to Worsley. Amer. Statist. 39 236.
\bibitem[281]{source281}\label{bib:source281} Schwarz, C. J. and Samanta, M. (1991). An Inductive Proof of the Sampling Distributions for the MLEs of the Parameters in an Inverse Gaussian Distribution. Amer. Statist. 45 223–225.
\bibitem[282]{source282}\label{bib:source282} Searle, S. R. (1971). Linear Models. New York: Wiley.
\bibitem[283]{source283}\label{bib:source283} Searle, S. R. (1982). Matrix Algebra Useful for Statistics. New York: Wiley.
\bibitem[284]{source284}\label{bib:source284} Searls, D. T. and Intarapanich, P. (1990). A Note on an Estimator for the Variance that Utilizes the Kurtosis. Amer. Statist. 44 295–296.
\bibitem[285]{source285}\label{bib:source285} Selvin, S. (1975). A Problem in Probability (letter to the editor). Amer. Statist. 29 67.
\bibitem[286]{source286}\label{bib:source286} Seshadri, V. (1993). The Inverse Gaussian Distribution. A Case Study in Exponential Families. New York: Clarendon Press.
\bibitem[287]{source287}\label{bib:source287} Shao, J. (1999). Mathematical Statistics. New York: Springer-Verlag.
\sourcepageanchor{552}
\bibitem[288]{source288}\label{bib:source288} Shao, J. and Tu, D. (1995). The Jackknife and the Bootstrap. New York: Springer-Verlag.
\bibitem[289]{source289}\label{bib:source289} Shier, D. P. (1988). The Monotonicity of Power Means Using Entropy. Amer. Statist. 42 203–204.
\bibitem[290]{source290}\label{bib:source290} Shuster, J. J. (1991). The Statistician in a Reverse Cocaine Sting. Amer. Statist. 45 123–124.
\bibitem[291]{source291}\label{bib:source291} Silvapulle, M. J. (1996). A Test in the Presence of Nuisance Parameters. J. Amer. Statist. Assoc. 91 1690–1693.
\bibitem[292]{source292}\label{bib:source292} Smith, A. F. M. and Gelfand, A. E. (1992). Bayesian Statistics Without Tears: A Sampling- Resampling Perspective. Amer. Statist. 46 84–88.
\bibitem[293]{source293}\label{bib:source293} Smith, A. F. M. and Roberts, G. O. (1993). Bayesian Computation via the Gibbs Sampler and Related Markov Chain Monte Carlo Methods (with discussion). J. Roy. Statist. Soc. Ser. B 55 3–24.
\bibitem[294]{source294}\label{bib:source294} Snedecor, G. W. and Cochran, W. G. (1989). Statistical Methods, 8th edition. Ames, IA: Iowa State University Press.
\bibitem[295]{source295}\label{bib:source295} Solari, M. E. (1969). The “Maximum Likelihood Solution” of the Problem of Estimating a Linear Functional Relationship. J. Roy. Statist. Soc. Ser. B 31 372–375.
\bibitem[296]{source296}\label{bib:source296} Solomon, D. L. (1975). A Note on the Non-Equivalence of the Neyman–Pearson and Generalized Likelihood Ratio Tests for Testing a Simple Null Versus a Simple Alternative Hypothesis. Amer. Statist. 29 101–102.
\bibitem[297]{source297}\label{bib:source297} Solomon, D. L. (1983). The Spatial Distribution of Cabbage Butterfly Eggs. Life Science Models, Volume 4 (H. Marcus-Roberts and M. Thompson, eds.) New York: Springer-Verlag 350–366.
\bibitem[298]{source298}\label{bib:source298} Sprott, D. A. and Farewell, V. T. (1993). The Difference Between Two Normal Means. Amer. Statist. 47 126–128.
\bibitem[299]{source299}\label{bib:source299} Staudte, R. G. and Sheather, S. J. (1990). Robust Estimation and Testing. New York: Wiley.
\bibitem[300]{source300}\label{bib:source300} Stefanski, L. A. (1996). A Note on the Arithmetic-Geometric-Harmonic Mean Inequalities. Amer. Statist. 50 246–247.
\bibitem[301]{source301}\label{bib:source301} Stein, C. (1956). Inadmissibility of the Usual Estimator for the Mean of a Multivariate Normal Distribution. Proc. Third Berkeley Symp. Math. Statist. Prob. 1. University of California Press 197–206.
\bibitem[302]{source302}\label{bib:source302} Stein, C. (1964). Inadmissibility of the Usual Estimator for the Variance of a Normal Distribution with Unknown Mean. Ann. Inst. Statist. Math. 16 155–160.
\bibitem[303]{source303}\label{bib:source303} Stein, C. (1973). Estimation of the Mean of a Multivariate Distribution. Proceedings of the Prague Symposium on Asymptotic Statistics. Prague: Charles Univ. 345–381.
\bibitem[304]{source304}\label{bib:source304} Stein, C. (1981). Estimation of the Mean of a Multivariate Normal Distribution. Ann. Statist. 9 1135–1151.
\bibitem[305]{source305}\label{bib:source305} Sterne, T. E. (1954). Some Remarks on Confidence or Fiducial Limits. Biometrika 41 275–278.
\bibitem[306]{source306}\label{bib:source306} Stigler, S. M. (1983). Who Discovered Bayes’ Theorem? Amer. Statist. 37 290–296.
\bibitem[307]{source307}\label{bib:source307} Stigler, S. M. (1984). Kruskal’s Proof of the Joint Distribution of $\bar X$ and $S^2$ . Amer. Statist. 38 134–135.
\bibitem[308]{source308}\label{bib:source308} Stigler, S. M. (1986). The History of Statistics: The Measurement of Uncertainty Before 1900. Cambridge, MA: Harvard University Press.
\bibitem[309]{source309}\label{bib:source309} Stuart, A. and Ord, J. K. (1987). Kendall’s Advanced Theory of Statistics, Volume I: Distribution Theory, 5th edition. New York: Oxford University Press.
\bibitem[310]{source310}\label{bib:source310} Stuart, A. and Ord, J. K. (1991). Kendall’s Advanced Theory of Statistics, Volume II, 5th edition. New York: Oxford University Press.
\sourcepageanchor{553}
\bibitem[311]{source311}\label{bib:source311} Stuart, A., Ord, J. K., and Arnold, S. (1999). Advanced Theory of Statistics, Volume 2A: Classical Inference and the Linear Model, 6th edition. London: Oxford University Press.
\bibitem[312]{source312}\label{bib:source312} Tanner, M. A. (1996) Tools for Statistical Inference: Observed Data and Data Augmentation Methods, 3rd edition. New York: Springer-Verlag.
\bibitem[313]{source313}\label{bib:source313} Tate, R. F. and Klett, G. W. (1959). Optimal Confidence Intervals for the Variance of a Normal Distribution. J. Amer. Statist. Assoc. 54 674–682.
\bibitem[314]{source314}\label{bib:source314} Tierney, J. (1991). Behind Monte Hall’s Door: Puzzle, Debate and Answer? The New York Times, July 21, 1991.
\bibitem[315]{source315}\label{bib:source315} Tierney, L. (1994). Markov Chains for Exploring Posterior Distributions (with discussion). Ann. Statist. 22 1701–1786.
\bibitem[316]{source316}\label{bib:source316} Tsao, C. A. and Hwang, J. T. (1998). Improved Confidence Estimators for Fieller’s Confidence Sets. Can. J. Statist. 26 299–310.
\bibitem[317]{source317}\label{bib:source317} Tsao, C. A. and Hwang, J. T. (1999). Generalized Bayes Confidence Estimators for Fieller’s Confidence Sets. Statistica Sinica 9 795–810.
\bibitem[318]{source318}\label{bib:source318} Tukey, J. W. (1977). Exploratory Data Analysis. Reading, MA: Addison-Wesley.
\bibitem[319]{source319}\label{bib:source319} Tweedie, M. C. K. (1957). Statistical Properties of Inverse Gaussian Distributions I. Ann. Math. Statist. 28 362–377.
\bibitem[320]{source320}\label{bib:source320} Vardeman, S. B. (1987). Discussion of the Articles by Casella and Berger and Berger and Sellke. J. Amer. Statist. Assoc. 82 130–131.
\bibitem[321]{source321}\label{bib:source321} Vardeman, S. B. (1992). What About Other Intervals. Amer. Statist. 46 193–197.
\bibitem[322]{source322}\label{bib:source322} vos Savant, M. (1990). Ask Marilyn. Parade Magazine, September 9, 15.
\bibitem[323]{source323}\label{bib:source323} vos Savant, M. (1991). Letter to the Editor. Amer. Statist. 45 347.
\bibitem[324]{source324}\label{bib:source324} Wald, A. (1940). The Fitting of Straight Lines when Both Variables Are Subject to Error. Ann. Math. Statist. 11 284–300.
\bibitem[325]{source325}\label{bib:source325} Waller, L. A. (1995). Does the Characteristic Function Numerically Distinguish Distributions? Amer. Statist. 49 150–151.
\bibitem[326]{source326}\label{bib:source326} Waller, L. A., Turnbull, B. W., and Hardin, J. M. (1995). Obtaining Distribution Functions by Numerical Inversion of Characteristic Functions with Applications. Amer. Statist. 49 346–350.
\bibitem[327]{source327}\label{bib:source327} Wasserman, L. (1992). Recent Methodological Advances in Robust Bayesian Inference (with discussion). Bayesian Statistics 4. J. M. Bernardo, J. O. Berger, A. P. Dawid and A. F. M. Smith, eds. Oxford: Clarendon Press, 483–502.
\bibitem[328]{source328}\label{bib:source328} Westlake, W. J. (1981). Bioequivalence Testing – A Need to Rethink. Biometrics 37 591–593.
\bibitem[329]{source329}\label{bib:source329} Widder, D. V. (1946). The Laplace Transform. Princeton, NJ: Princeton University Press.
\bibitem[330]{source330}\label{bib:source330} Wilks, S. S. (1938). Shortest Average Confidence Intervals from Large Samples. Ann. Math. Statist. 9 166–175.
\bibitem[331]{source331}\label{bib:source331} Williams, E. J. (1959). Regression Analysis. New York: Wiley.
\bibitem[332]{source332}\label{bib:source332} Worsley, K. J. (1982). An Improved Bonferroni Inequality and Applications. Biometrika 69 297–302.
\bibitem[333]{source333}\label{bib:source333} Worsley, K. J. (1985). Bonferroni (Improved) Wins Again. Amer. Statist. 39 235.
\bibitem[334]{source334}\label{bib:source334} Wright, T. (1992). Lagrange’s Identity Reveals Correlation Coefficient and Straight-Line Connection. Amer. Statist. 46 106–107.
\bibitem[335]{source335}\label{bib:source335} Wu, C. F. J. (1983). On the Convergence of the EM Algorithm. Ann. Statist. 11 95–103.
\sourcepageanchor{554}
\bibitem[336]{source336}\label{bib:source336} Young, G. A. (1994). Bootstrap: More Than a Stab in the Dark? (with discussion). Statist. Sci. 9 382–415.
\bibitem[337]{source337}\label{bib:source337} Zehna, P. W. (1966). Invariance of Maximum Likelihood Estimators. Ann. Math. Statist. 37 744.
\bibitem[338]{source338}\label{bib:source338} Zellner, A. (1986). Bayesian Estimation and Prediction Using Asymmetric Loss Functions. J. Amer. Statist. Assoc. 81 446–451.
\end{thebibliography}
